\UseRawInputEncoding
\chapter{Calcolo combinatorio}


\begin{esercizio}
Quanti sono i numeri naturali formati da quattro cifre tutte distinte tra loro?
\end{esercizio}
\textit{Soluzione.} Occorre contare quanti numeri iniziano con una qualsiasi cifra diversa da zero e non hanno cifre ripetute:
\[
N=9\cdot 9\cdot 8\cdot 7=9\cdot D_{9,3}
\]

\begin{esercizio}
	Quante colonne si devono compilare per essere certi di fare $13$ al totocalcio?
\end{esercizio}
\textit{Soluzione.} Le $13$ colonne vengono compilate inserendo, con ripetizione, i tre simboli disponibili. Il numero di colonne � quindi uguale al numero di disposizioni con ripetizione:  
\[
D'_{13,3}=3^{13}
\] 

\begin{esercizio}
	Quanti interi compresi fra $100$ e $1000$ si possono comporre con le cifre: $2,4,7,8$?
\end{esercizio}
\textit{Soluzione.} Si osserva che i numeri richiesti sono formati da tre cifre, anche ripetute, e che non si pone il problema di escludere lo zero per la cifra delle centinaia: 
\[
D'_{4,3}=4^3
\]

\begin{esercizio}
	Quante combinazioni si possono ottenere nel lancio di due dati?
\end{esercizio}
\textit{Soluzione.} Si tratta del numero di combinazioni con ripetizioni: 
\[
C'_{6,2}=C_{6+1, 2}={7\choose 2}=21
\]

\begin{esercizio}
	Quanti anagrammi si possono formare rispettivamente con le parole \textit{studiare} e con \textit{esponente}?
\end{esercizio}
\textit{Soluzione.} Le lettere della parola \textit{studiare} non sono ripetute. Il numero di anagrammi � uguale a quello delle permutazioni semplici: 
\[
P_{8}=8!
\]
Nella parola \textit{anagrammi}, alcune lettere si ripetono. Occorre quindi considerare il numero di permutazioni con ripetizione. 
\[
P'_{3,2,1,1,1,1}=\frac{9!}{3! \, 2!}
\]

\begin{esercizio}
	Sviluppare il seguente binomio:
	\[
	\Big (\sqrt{x}-\frac{1}{2} \, y\Big )^5
	\]
\end{esercizio}
\textit{Soluzione.} 
\begin{align*}
	\Big (\sqrt{x}-\frac{1}{2}\, y\Big )^5
	&=\sum_{k=0}^5 {5\choose k} \big (\sqrt{x}\big )^k \cdot \Big (\frac{-y}{2}\Big )^{5-k}\\
	&=x^2\sqrt{x}-\frac{5}{2} \, x^2y+\frac{5}{2} \, x \sqrt{x} \, y^2-\frac{5}{4} \, xy^3+\frac{5}{16} \, \sqrt{x} \, y^4-\frac{1}{32} \, y^5
\end{align*}

\begin{esercizio}
	Calcolare il coefficiente del termine in $x^2y^2$ nel seguente binomio: 
	\[
	\Big (\frac{y}{x}+\sqrt{x}\Big )^{10}
	\]
\end{esercizio}
\textit{Soluzione.} Il termine $k$-esimo risulta:
\[
{10\choose k} \, \Big (\frac{y}{x}\Big )^k \, \big (\sqrt{x}\big )^{10-k}
\quad \Rightarrow \quad 
{10 \choose k} \, \frac{y^k}{x^k} \cdot x^{(10-k)/2}=
{10 \choose k} \, y^k \cdot x^{(10-3k)/2}
\]
Confrontando le potenze di $x$ e di $y$:
\[
\begin{cases}
	\dfrac{10-3k}{2}=2 \\[6pt]
	k=2	
\end{cases}
\Rightarrow \quad 
\begin{cases}
	10-6=4
	\\
	k=2
\end{cases}
\Rightarrow \quad 
k=2
\]

\begin{esercizio}
	Calcolare il coefficiente del termine indipendente da $x$ del termine in $x^2y^2$ nel binomio: 
	\[
	\Big (\frac{y}{x}+\sqrt{x}\Big )^{10}
	\]
\end{esercizio}
\textit{Soluzione.} Il termine $k$-esimo risulta:
\[
{10\choose k} \, \Big (\frac{y}{x}\Big )^k \, \big (\sqrt{x}\big )^{10-k}
\quad \Rightarrow \quad 
{10 \choose k} \, \frac{y^k}{x^k} \cdot x^{(10-k)/2}=
{10 \choose k} \, y^k \cdot x^{(10-3k)/2}
\]
Il termine indipendente da $x$ soddisfa la relazione: 
\[
\frac{10-3k}{2}=0 \quad \Rightarrow \quad 
k=\frac{10}{3}
\]
Il valore di $k$ non � accettabile in quanto $k$ deve essere intero. 

\begin{esercizio}
	Dimostrare le seguenti uguaglianze:
	\[
		\sum_{k=0}^n {n\choose k}=2^n
		\qquad \quad 
		\sum_{k=0}^n (-1)^k {n\choose k}=0
	\]
\end{esercizio}
\textit{Soluzione.} 
\begin{itemize}
	\item Si applica il binomio di Newton ponendo $a=b=1$:
	\[
	(1+1)^n=\sum_{k=0}^n {n\choose k} 1^k \cdot 1^{n-k}=\sum_{k=0}^n {n\choose k} 	
	\quad \Rightarrow \quad 
	\sum_{k=0}^n {n\choose k}=2^n
	\]
	\item Si applica il binomio di Newton ponendo $a=1$ e $b=-1$:
	\[
	\big (1+(-1)\big )^n=\sum_{k=0}^n {n\choose k} (-1)^k \cdot 1^{n-k}=\sum_{k=0}^n {n\choose k} (-1)^k
	\quad \Rightarrow \quad 
	\sum_{k=0}^n (-1)^k \, {n\choose k}=0 
	\]
	Dallo sviluppo precedente si deduce che la somma dei coefficienti di indice pari � uguale allo somma di quelli con indice dispari: 
	\[
	{n\choose 0}+{n\choose 2}+{n\choose 4}+\dots
	={n\choose 1}+{n\choose 3}+{n\choose 5}+\dots
	\]
	Da cui: 
	\[
	{n\choose 0}+{n\choose 2}+{n\choose 4}+\dots
	={n\choose 1}+{n\choose 3}+{n\choose 5}+\dots
	=\frac{2^n}{2}=2^{n-1}
	\]
\end{itemize}

% Novit� GIT